In mutualised logistics networks the urgency an operator declares often diverges from the
urgency computable from the operational state. Over-declaration wastes shared capacity;
under-declaration hides rupture risk. No existing digital twin measures that divergence or
propagates it across a supplier network.

This thesis builds and tests one. Real urgency is aggregated at each node from six KPI blocks
through a weighted probabilistic OR and propagated upstream across a directed acyclic supply
graph; declared urgency is elicited through a consistency-gated pairwise questionnaire and
propagated downstream; their mismatch is scored on $[0,100]$ with an asymmetric
Prospect-Theory penalty. The system was evaluated through a pre-registered campaign replaying
the 2020 to 2022 semiconductor shortage over nineteen turns and eight nodes, with a blind
control arm and a forecast-assisted treatment arm on a frozen scenario.

The results are negative and diagnostically precise. The forecast layer shows no
discrimination above chance and is severely miscalibrated, at a Brier skill between $-3.4$
and $-6.9$ against a base-rate forecaster. Forecasts below $0.1$ are well calibrated, while
twenty-one announcing a mean probability of $0.948$ were followed by the predicted outcome
exactly zero times. One block explains it: the time block reads a lead-time distribution that
no shock ever reaches. In the treatment arm, on the turn the forecast first became visible,
all eight declarants revised their urgency downward or held it and none revised upward, two
turns before the scenario's critical event. Both arms carry identical forecasts, so the effect
belongs to the display. Substituting a plausible but naive outcome definition improves the
apparent skill fivefold while leaving discrimination unchanged.

\vspace{0.5cm}

\noindent\textbf{Keywords:} supply chain digital twin, urgency adequacy, cognitive bias,
multi-criteria decision analysis, forecast calibration, automation bias, serious game.
